% Part: turing-machines
% Chapter: machines-computations
% Section: combining-machines

\documentclass[../../../include/open-logic-section]{subfiles}

\begin{document}

\olfileid{tur}{mac}{cmb}
\olsection{组合Turing机}

\begin{explain}
迄今为止我们所见的图灵机示例都相当简单。
但事实上，任何能用现代编程语言解决的问题，同样也能用图灵机解决。
要构建更复杂的图灵机，关键在于确信我们能够将它们组合起来，从而可以通过将计算过程分解为更简单的部分，来构建解决更复杂问题的机器。
如果我们能找到一种自然的方式将复杂问题分解为若干组成部分，
就可以分阶段处理该问题：先构建几台简单的图灵机，再将它们组合成一台能够解决该问题的机器。
这一点在下一节处理停机问题时尤为重要。

我们如何组合图灵机 $M = \tuple{Q, \Sigma, q_0, \delta}$ 与 $M' = \tuple{Q', \Sigma', q_0', \delta'}$ 呢？
现在，我们将 $M$ 停机后的纸带配置用作机器 $M'$ 运行的输入配置。
为了得到一台能完成此事的单一图灵机 $M \frown M'$，需执行以下操作：
\begin{enumerate}
    \item 对 $M'$ 的所有状态 $Q'$ 重新编号（或重新标记），使得 $M$ 和 $M'$ 没有公共状态（$Q \cap Q' = \emptyset$）。
    \item $M \frown M'$ 的状态集为 $Q \cup Q'$。
    \item 纸带字母表为 $\Sigma \cup \Sigma'$。
    \item 起始状态为 $q_0$。
    \item 转移函数是如下定义的函数 $\delta''$：
    \[\delta''(q,\sigma) =
    \begin{cases}
      \delta(q,\sigma) & \text{若 $q \in Q$}\\
      \delta'(q,\sigma) & \text{若 $q \in Q'$}\\
      \tuple{q_0', \sigma, \TMstay} & \text{若 $q \in Q$ 且
      $\delta(q,\sigma)$ 无定义}
    \end{cases}\]
\end{enumerate}
这样得到的机器，当处于状态 $q \in Q$ 时使用 $M$ 的指令，当处于状态 $q \in
Q'$ 时使用 $M'$ 的指令。当它处于状态 $q \in Q$ 且正在扫描符号 $\sigma$，
而 $M$ 对该符号没有转移（即 $M$ 本应停机）时，它就进入 $M'$ 的起始状态（并保持纸带内容和读写头位置不变）。

注意，除非机器 $M$ 是规整的，否则当 $M$ 停机时，我们不知道读写头在什么位置，
因此 $M$ 的停机配置未必能让读写头扫描方格 $1$。
在组合机器时，务必牢记这一点。
\end{explain}

\begin{ex}
\emph{机器组合：}我们将设计一台机器，当输入由两个长度分别为 $n$ 和 $m$ 的 $\TMstroke$ 划痕块组成时，它停机并在纸带上留下一个长度为 $2(m+n)$ 的 $\TMstroke$ 划痕块。为了构建这台机器，我们可以组合两台我们已经熟悉的机器：加法机和加倍机。我们首先画出加法机的状态图。
\[
\begin{tikzpicture}[->,>=stealth',shorten >=1pt,auto,node distance=2.8cm,
                    semithick]
  \tikzstyle{every state}=[fill=none,draw=black,text=black]

  \node[initial,state] (A)              {$q_0$};
  \node[state]         (B) [right of=A] {$q_1$};
  \node[state]         (C) [right of=B] {$q_2$};

  \path (A) edge node {\TMtrans{\TMblank}{\TMstroke}{\TMstay}} (B)
            edge [loop above] node {\TMtrans{\TMstroke}{\TMstroke}{\TMright}} (B)
        (B) edge [loop above] node {\TMtrans{\TMstroke}{\TMstroke}{\TMright}} (B)
            edge node {\TMtrans{\TMblank}{\TMblank}{\TMleft}} (C)
        (C) edge [loop above] node {\TMtrans{\TMstroke}{\TMblank}{\TMstay}} (C);
\end{tikzpicture}
\]
我们不希望在状态 $q_2$ 停机，而是希望继续运行以将输出加倍。回想一下，加倍机会擦除输入中的第一个划痕，并在单独的输出区写下两个划痕。让我们添加一条指令，以确保读写头正读取加法机输出的第一个划痕。
\[
\begin{tikzpicture}[->,>=stealth',shorten >=1pt,auto,node distance=2.8cm,
                    semithick]
  \tikzstyle{every state}=[fill=none,draw=black,text=black]

  \node[initial,state] (A)              {$q_0$};
  \node[state]         (B) [right of=A] {$q_1$};
  \node[state]         (C) [right of=B] {$q_2$};
  \node[state]         (D) [below left of=C] {$q_3$};
  \node[state]         (E) [below left of=D] {$q_4$};

  \path (A) edge node {\TMtrans{\TMblank}{\TMstroke}{\TMstay}} (B)
            edge [loop above] node {\TMtrans{\TMstroke}{\TMstroke}{\TMright}} (B)
        (B) edge [loop above] node {\TMtrans{\TMstroke}{\TMstroke}{\TMright}} (B)
            edge node {\TMtrans{\TMblank}{\TMblank}{\TMleft}} (C)
        (C) edge node {\TMtrans{\TMstroke}{\TMblank}{\TMleft}} (D)
        (D) edge [loop left] node {\TMtrans{\TMstroke}{\TMstroke}{\TMleft}} (D)
            edge node[left, xshift=-2mm] {\TMtrans{\TMendtape}{\TMendtape}{\TMright}} (E);
\end{tikzpicture}
\]
现在要使输入加倍就很容易了——我们只需将加倍机连接到状态 $q_4$ 上。这需要重命名加倍机的状态，使其从 $q_4$ 开始而非 $q_0$——这样我们才不会出现两个起始状态。最终的图应如 \olref{fig:combined} 所示。
\begin{figure}
\[
\begin{tikzpicture}[->,>=stealth',shorten >=1pt,auto,node distance=2.8cm,
                    semithick]
  \tikzstyle{every state}=[fill=none,draw=black,text=black]
  \node[initial,state] (A)              {$q_0$};
  \node[state]         (B) [right of=A] {$q_1$};
  \node[state]         (C) [right of=B] {$q_2$};
  \node[state]         (D) [below left of=C] {$q_3$};
  \node[state]         (E) [below left of=D] {$q_4$};
  \node[state]         (2) [right of=E] {$q_5$};
  \node[state]         (3) [right of=2] {$q_6$};
  \node[state]         (4) [below of=3] {$q_7$};
  \node[state]         (5) [left of=4]  {$q_8$};
  \node[state]         (6) [left of=5]  {$q_9$};

  \path (A) edge node {\TMtrans{\TMblank}{\TMstroke}{\TMstay}} (B)
            edge [loop above] node {\TMtrans{\TMstroke}{\TMstroke}{\TMright}} (B)
        (B) edge [loop above] node {\TMtrans{\TMstroke}{\TMstroke}{\TMright}} (B)
            edge node {\TMtrans{\TMblank}{\TMblank}{\TMleft}} (C)
        (C) edge node {\TMtrans{\TMstroke}{\TMblank}{\TMleft}} (D)
        (D) edge [loop left] node {\TMtrans{\TMstroke}{\TMstroke}{\TMleft}} (D)
            edge node[left, xshift=-2mm] {\TMtrans{\TMendtape}{\TMendtape}{\TMright}} (E)
    (E) edge node {\TMtrans{\TMstroke}{\TMblank}{\TMright}} (2)
    (2) edge [loop above] node {\TMtrans{\TMstroke}{\TMstroke}{\TMright}} (2)
      edge node {\TMtrans{\TMblank}{\TMblank}{\TMright}} (3)
    (3) edge [loop above] node {\TMtrans{\TMstroke}{\TMstroke}{\TMright}} (3)
        edge node {\TMtrans{\TMblank}{\TMstroke}{\TMright}} (4)
    (4) edge [loop below] node {\TMtrans{\TMblank}{\TMstroke}{\TMleft}} (4)
        edge node {\TMtrans{\TMstroke}{\TMstroke}{\TMleft}} (5)
    (5) edge [loop below]  node {\TMtrans{\TMstroke}{\TMstroke}{\TMleft}} (5)
        edge              node {\TMtrans{\TMblank}{\TMblank}{\TMleft}} (6)
    (6) edge [loop below] node {\TMtrans{\TMstroke}{\TMstroke}{\TMleft}} (6)
        edge              node {\TMtrans{\TMblank}{\TMblank}{\TMright}} (E);
\end{tikzpicture}
\]
\caption{组合加法机与加倍机}
\ollabel{fig:combined}
\end{figure}
\end{ex}

\begin{prop}
    若 $M$ 和 $M'$ 是规整的，且分别计算函数 $f\colon
    \Nat^k \to \Nat$ 和 $f'\colon \Nat \to \Nat$，则 $M \frown M'$ 是规整的且计算~$\comp{f}{f'}$。
\end{prop}

\begin{proof}
    由于 $M$ 是规整的，当它带着输出 $f(n_1,\dots,n_k) = m$ 停机时，读写头正扫描方格 $1$。若我们此时进入 $M'$ 的起始状态，机器将带着输出 $f'(m)$ 停机，并再次扫描方格 $1$。同样满足 \olref[dis]{defn:disciplined} 的其他条件。
\end{proof}

\begin{prob}
    取 \cref{tur:mac:dis:prob:disc-succ} 中的机器 $M$，构造 $M \frown M$，从而给出一台计算 $f(x) = x+2$ 的规整图灵机。
\end{prob}
\end{document}